\documentclass[../../../include/open-logic-section]{subfiles}

\begin{document}
\olfileid{sth}{card-arithmetic}{simp}

\olsection{Nyederhanakaké Panambahan lan Ping-pingan}

Ternyata panambahan lan ping-pingan kardinal transfinit iku
\emph{gampang banget}. Iki amarga kardinal iku ordinal
(tartamtu), mula terurut becik lan bisa diolah kanthi cara
tartamtu. Nanging mbuktèkaké iki \emph{ora} gampang.
Kanggo miwiti, kita perlu definisi kang rada njlimet:

\begin{defn}
Kita netepaké \emph{pengurutan kanonik}, $\canonord$,
ing pasangan ordinal, kanthi aturan yèn $\tuple{\alpha_1, \alpha_2} \canonord
\tuple{\beta_1, \beta_2}$ yen lan mung yen salah siji:
\begin{enumerate}
	\item $\max(\alpha_1, \alpha_2) < \max(\beta_1, \beta_2)$; utawa
	\item $\max(\alpha_1, \alpha_2) = \max(\beta_1, \beta_2)$ lan
	$\alpha_1 < \beta_1$; utawa
	\item $\max(\alpha_1, \alpha_2) = \max(\beta_1, \beta_2)$ lan
	$\alpha_1 = \beta_1$ lan $\alpha_2 < \beta_2$
\end{enumerate}
\end{defn}

\begin{lem}
$\tuple{\alpha \times \alpha, \canonord}$ iku pengurutan becik,
kanggo sembarang ordinal $\alpha$.
\end{lem}

\begin{proof}
Cetha yèn $\canonord$ nyambung ing $\alpha \times \alpha$.
Upamané $\tuple{\alpha_1, \alpha_2}$ lan $\tuple{\beta_1,
\beta_2}$ ora ana kang luwih cilik miturut $\canonord$
tinimbang sijiné. Mula $\max(\alpha_1,
\alpha_2) = \max(\beta_1, \beta_2)$ lan $\alpha_1 = \beta_1$ lan
$\alpha_2 = \beta_2$, saéngga $\tuple{\alpha_1, \alpha_2} =
\tuple{\beta_1,  \beta_2}$.

Kanggo nuduhaké sipat pengurutan becik, tetepna
$X \subseteq \alpha\times\alpha$ kang ora kosong.
Amarga $\alpha$ ordinal, ana $\delta$ kang dadi unsur
paling cilik ing $\Setabs{\max(\gamma_1, \gamma_2)}{\tuple{\gamma_1,
\gamma_2} \in X}$. Saiki sisihna kabèh pasangan saka
$\Setabs{\tuple{\gamma_1,\gamma_2} \in X}{\max(\gamma_1, \gamma_2) =
\delta}$ kajaba kang koordinat kapisané paling cilik;
ing antarané, pasangan kang koordinat kapindhoné
paling cilik iku unsur paling cilik miturut $\canonord$ ing $X$.
\end{proof}
\noindent
Saiki sawijining pangamatan prasaja:

\begin{prop}\ollabel{simplecardproduct}
Yèn $\cardeq{\alpha}{\beta}$, mula $\cardeq{\alpha \times \alpha}{\beta
\times \beta}$.
\end{prop}

\begin{proof}
Cukup gunakna $f \colon \alpha \to \beta$ kanggo mbangun
pemetaan $\tuple{\gamma_1,
\gamma_2} \mapsto \tuple{f(\gamma_1), f(\gamma_2)}$.
\end{proof}
\noindent
Saiki kabèh iki bakal dienggo kanggo mbuktèkaké lema wigati:
\begin{lem}\ollabel{alphatimesalpha}
$\cardeq{\alpha}{\alpha \times \alpha}$ kanggo sembarang
ordinal tanpa wates $\alpha$.
\end{lem}

\begin{proof}
Kanggo kontradiksi, upamané $\alpha$ ordinal tanpa wates paling
cilik kang ora nyukupi iki.
\olref[sfr][siz][zigzag]{natsquaredenumerable} nuduhaké
$\cardeq{\omega}{\omega\times\omega}$, mula $\omega \in \alpha$.
Kajaba iku, $\alpha$ iku kardinal: yèn ora, kanggo
kontradiksi, $\card{\alpha} \in \alpha$, saéngga
$\cardeq{\card{\alpha}}{\card{\alpha} \times \card{\alpha}}$ miturut
hipotesis; lan $\cardeq{\card{\alpha}}{\alpha}$ miturut
definisi; mula $\cardeq{\alpha}{\alpha\times\alpha}$ miturut
\olref{simplecardproduct}.

Saiki, kanggo saben $\tuple{\gamma_1, \gamma_2} \in \alpha \times \alpha$,
gatekna segmen iki:
\begin{align*}
	\text{Seg}(\gamma_1, \gamma_2) &= \Setabs{\tuple{\delta_1, \delta_2} \in \alpha \times \alpha}{\tuple{\delta_1, \delta_2} \canonord \tuple{\gamma_1, \gamma_2}}
\end{align*}
Tetepna $\gamma = \max(\gamma_1, \gamma_2)$; wigatèkna yèn $\tuple{\gamma_1, \gamma_2} \canonord \tuple{\gamma+1, \gamma + 1}$. Dadi, nalika $\gamma$ tanpa wates, kita olèh:
\begin{align*}
	\text{Seg}(\gamma_1, \gamma_2) &
	\precsim ((\gamma \ordplus 1)\ordtimes (\gamma \ordplus 1))\\
	&\approx (\gamma \ordtimes \gamma)
	\text{, miturut \olref[ord-arithmetic][using-addition]{ordinfinitycharacter} lan
	\olref{simplecardproduct}}\\
	&\approx \gamma \text{, miturut hipotesis induksi}\\
	& \prec \alpha\text{, amarga $\alpha$ iku kardinal}
\end{align*}
Yèn indeks maksimum mau winates, segmen kasebut uga winates, mula cacahé
luwih cilik tinimbang kardinal tanpa wates iki. Dadi
$\ordtype{\alpha\times \alpha, \canonord} \leq \alpha$, mula
$\cardle{\alpha \times \alpha}{\alpha}$. Amarga cetha
$\cardle{\alpha}{\alpha \times \alpha}$, asilé nututi
teorema Schr\"oder-Bernstein.
\end{proof}

Pungkasané, kita tekan asil panyederhanaan:

\begin{thm}\ollabel{cardplustimesmax}
Yèn $\cardfont{a}, \cardfont{b}$ kardinal tanpa wates, mula:
\[
	\cardfont{a}
\cardtimes \cardfont{b} = \cardfont{a} \cardplus \cardfont{b} =
\text{max}(\cardfont{a}, \cardfont{b}).
\]
\end{thm}

\begin{proof}
Tanpa ngurangi keumuman, upamané $\cardfont{a} = \max(\cardfont{a},
\cardfont{b})$. Kanthi \olref{alphatimesalpha},
$\cardfont{a}\cardtimes\cardfont{a} = \cardfont{a} \leq \cardfont{a}
\cardplus \cardfont{b} \leq \cardfont{a} \cardplus \cardfont{a} \leq
\cardfont{a} \cardtimes \cardfont{a}$. \end{proof}\noindent Semono uga,
yèn $\cardfont{a}$ tanpa wates, gabungan saka
$\cardfont{a}$ himpunan kang saben-saben ukurané
$\leq\cardfont{a}$ uga nduwèni ukuran $\leq\cardfont{a}$:

\begin{prop}\ollabel{kappaunionkappasize}
Upamané $\cardfont{a}$ kardinal tanpa wates. Kanggo saben ordinal $\beta
\in \cardfont{a}$, upamané $X_\beta$ himpunan kang $\card{X_\beta} \leq
\cardfont{a}$. Mula $\card{\bigcup_{\beta \in \cardfont{a}} X_\beta}
\leq \cardfont{a}$.
\end{prop}

\begin{proof}
Kanggo saben $\beta \in \cardfont{a}$, pilih !!a{injection} $f_\beta \colon
X_\beta \to \cardfont{a}$.\footnote{Kepiyé pilihan-pilihan iki ditetepaké? Delengen \olref[sth][choice][countablechoice]{sec}.} Tetepna !!a{injection} $g \colon
\bigcup_{\beta \in \cardfont{a}} X_\beta \to \cardfont{a} \times
\cardfont{a}$ kanthi $g(v) = \tuple{\beta, f_\beta(v)}$, ing kéné $v \in
X_\beta$ lan $v \notin X_\gamma$ kanggo sembarang $\gamma \in \beta$. Saiki
$\bigcup_{\beta \in \cardfont{a}} X_\beta \preceq \cardfont{a} \times
\cardfont{a} \approx \cardfont{a}$ miturut \olref{cardplustimesmax}.
\end{proof}

\end{document}
